% !TeX program = pdflatex
% !TeX encoding = UTF-8
% !BIB program = biber
% !TeX root = main.tex
% 正文需要输出中文时：改用 XeLaTeX（本行改为 xelatex，并对调 .vscode/settings.json 中的注释）
%
% 说明：本骨架是一份"正常的 Stage 1 进展报告"——按研究本身的逻辑组织章节，
% 只在内容上自然覆盖官方 Stage 1 的六条标准（背景与批判性认识、研究方法、
% 批判分析与贡献、学术交流、计划与风险、培训与发展），不另设评分对照表。
\documentclass{article}
\usepackage{Temp/Report_Temp}

\usepackage{tabularx}
\usepackage{enumitem}
\usepackage{array}

% 表格列：左对齐、不拉伸、避免难看的连字符断词
\newcolumntype{L}[1]{>{\raggedright\arraybackslash}p{#1}}
\newcolumntype{Y}{>{\raggedright\arraybackslash}X}

% Choose the refer style here
% \usepackage{Temp/Harvard_Temp}
\usepackage{Temp/IEEE_Temp}

\begin{document}

\begin{titlepage} % The cover page, The only thing need to be change is the contents
    \begin{center}
        \thispagestyle{empty} % Clear the page numbers
        \vspace*{-3cm}
        \makebox[\textwidth]{\includegraphics[width=\paperwidth]{image/unnc_logo.jpg}}
        \vspace{0.2cm} % the distance between the logo and the school name

        {\bfseries\fontsize{18pt}{20pt}\selectfont [Your School / Department]}\\
        \vspace{0.8cm}
        {\bfseries\fontsize{16pt}{18pt}\selectfont [Your Lab / Research Group]}\\
        \vspace{1.2cm}

        {\bfseries\fontsize{18pt}{20pt}\selectfont PhD Progression Review Report}\\
        \vspace{0.4cm}
        {\bfseries\fontsize{15pt}{17pt}\selectfont Stage 1 --- Confirmation of PhD Registration}\\
        \vspace{1.0cm}
        {\bfseries\fontsize{15pt}{17pt}\selectfont Academic Year [20XX/XX]}\\
        \vspace{2.0cm}

        {\bfseries\fontsize{18pt}{20pt}\selectfont [First Line of Your Report Title]}\\
        \vspace{0.8cm}
        {\bfseries\fontsize{18pt}{20pt}\selectfont [Second Line of Your Report Title]}\\
        \vspace{2.8cm}

         \begin{flushleft}
            \begin{tabular}{l l}
                Candidate: & [Your Name]\\
                Student ID: & [Your Student ID]\\
                Supervisor(s): & [Dr Supervisor Name] \\
                Internal Assessor: & [Dr Internal Assessor] \\
                Review Period: & [MM/YYYY] -- [MM/YYYY] \\
                Date: & [DD/MM/YYYY]\\
            \end{tabular}
        \end{flushleft}
    \end{center}
\end{titlepage}

\tableofcontents

\newpage
% ============ 摘要：一页以内，让评审人先知道"做了什么、做到哪、接下来怎么办" ============
\section*{Abstract}
\addcontentsline{toc}{section}{Abstract}

This report reviews the first year of the doctoral project [short title of the project]. The
work is motivated by [the practical or theoretical problem the project addresses] and aims to
[the overall aim, in one sentence]. During the review period, [what has been done: e.g. the
literature review, the system model, the experimental platform] has been completed, and
[the main preliminary finding] has been obtained. The remaining work is organised into [number]
work packages, leading to [the expected contribution] and to thesis submission in month
[36]. The report also describes the training completed, the supervisory meetings held, and the
difficulties encountered during the year.

% ============ 1. 引言：背景、缺口、目标、范围 ============
\section{Introduction}

\subsection{Background and Motivation}
% 对应 Stage 1 标准 1：对领域现有知识的理解、对当前问题的批判性认识、研究理据。
The starting point of this project is [the technical or societal need that motivates the
research]. Over the past decade, [the field] has developed rapidly: [the key advances], and
[the enabling technologies] have made [what used to be impossible] feasible. In parallel, [the application domain] has placed new demands on
[performance, cost, reliability, efficiency], and these demands are difficult to meet with
existing solutions.

[Name the current dominant approach] is the most widely adopted solution at present. It works
well when [conditions], but its performance degrades sharply when [the conditions that do not
hold in practice] \cite{vaswani2017attention}. This mismatch between the assumptions behind
existing methods and the conditions found in real deployments is the practical problem that
this project addresses.

The motivation for the project is therefore twofold. From a research perspective, [the
scientific question that remains open]. From an application perspective, [who would benefit
and how]. Both perspectives shape the objectives set out in Section 1.3.

\subsection{Research Gap}
Despite the progress described above, three issues remain unresolved. First, [gap 1: what
existing work does not model, measure or explain]. Second, [gap 2: the limitation that follows
from the assumptions or the experimental conditions of previous studies]. Third, [gap 3: the
absence of a method that connects the two]. These gaps are discussed in more detail in
Section 2.3.

The gap that this project targets is [the specific gap, stated in one sentence]. Closing it
matters because [the consequence for theory or practice], and it is tractable because [the
recent development that now makes it possible].

\subsection{Aims, Research Questions and Objectives}
% 目标、问题、任务三层要能互相对上：问题 -> 工作包 -> 交付物。
The overall aim of the project is to [develop / derive / validate / quantify] [the object of
the research] so that [the improvement or understanding that is sought].

This aim is pursued through three research questions:
\begin{enumerate}[leftmargin=*]
    \item \textbf{RQ1.} [Question about the underlying mechanism, model or behaviour.]
    \item \textbf{RQ2.} [Question about the method, algorithm or design.]
    \item \textbf{RQ3.} [Question about validation, trade-offs or generalisation.]
\end{enumerate}

The questions translate into the following objectives, each associated with a work package:
\begin{enumerate}[leftmargin=*]
    \item \textbf{O1 (WP1).} [Objective, phrased with a verifiable verb, and the criterion by
    which completion will be judged.]
    \item \textbf{O2 (WP2).} [Objective and completion criterion.]
    \item \textbf{O3 (WP3).} [Objective and completion criterion.]
\end{enumerate}

\subsection{Scope and Assumptions}
The study is limited to [the operating conditions, the class of systems, the data available].
It does not attempt to [what is deliberately excluded], because [the reason]. The plan rests on
three assumptions: [assumption 1], [assumption 2] and [assumption 3]. If [assumption 2] does not
hold, [the alternative route that would be taken]; this is reflected in the risk assessment in
Section 5.3.

\subsection{Report Structure}
Section 2 reviews the literature and identifies the gap. Section 3 describes the methodology.
Section 4 reports the work completed and the preliminary results. Section 5 sets out the plan
for the remaining period, including the risk assessment and the thesis timetable. Section 6
covers training, academic engagement and supervision, Section 7 discusses the difficulties
encountered, and Section 8 concludes.

% ============ 2. 文献综述：不只是罗列，要有评价和缺口 ============
\section{Literature Review}

\subsection{Theoretical Background}
% 交代理解本研究所需的基础概念与模型，篇幅不宜长。
[Introduce the concepts and models on which the project builds: the governing theory, the
standard formulation, the metrics used in the field. Cite the foundational references.]

\subsection{State of the Art}
% 按主题组织，每个主题讲清"主流做法 + 达到的水平 + 遗留问题"。
The literature relevant to this project falls into three strands. The first strand concerns
[theme 1]. Representative work includes [author, year], who [approach and result], and
[author, year], who [approach and result]. The second strand addresses [theme 2] through
[methods], and has established [what is now well understood]. The third strand, closer to this
project, combines [theme 1] and [theme 2] in order to [purpose]; examples include
\cite{vaswani2017attention, lecun2015deep}.

Across these strands, [the consensus that has been reached] is well supported. What remains
less clear is [the point on which the strands disagree or fall silent].

\subsection{Critical Assessment and Research Gap}
% 这里是"批判性认识"的落点：评价证据强度，而不是复述结论。
The existing evidence is uneven. Studies based on [method A] report [results], but they rely on
[assumption] and rarely test [condition], so their conclusions may not transfer to [the setting
of this project]. Studies based on [method B] avoid that assumption but pay for it with
[limitation]. Only a small number of works consider both, and those that do
[e.g. use a simplified model, or report a single operating point].

This assessment points to the gap introduced in Section 1.2: [restate the gap in one sentence,
now supported by the literature]. The project addresses it by [the specific angle], which
distinguishes the present work from [the closest existing study].

\subsection{Implications for This Project}
The review has three consequences for the design of the study. It confirms that [design choice
1] is appropriate; it suggests that [design choice 2] should be preferred over the alternatives;
and it identifies [the benchmark or baseline] against which the results must be compared in
order to be convincing.

% ============ 3. 方法：设计、关键技术、数据与设备、局限 ============
\section{Research Methodology}

\subsection{Overall Research Design}
% 说明"用什么证据回答哪个问题"，这是方法章的主线。
The project follows a [theoretical / simulation-based / experimental / mixed] design, organised
around the three work packages introduced in Section 1.3. WP1 addresses RQ1 by [approach] and
produces [the model or dataset]. WP2 addresses RQ2 by [approach] and produces [the method or
design]. WP3 addresses RQ3 by [approach] and produces [the validation evidence]. The logic is
that [why this sequence of evidence is sufficient to answer the questions].

\subsection{Key Methods and Techniques}
% 对应标准 2：既要说清"用了哪些方法"，也要说清"为什么是这些方法、掌握到什么程度"。
The methods listed in Table~\ref{tab:methods} form the technical core of the project. Each was
selected after comparing it with the alternatives available in the literature; the third column
records the reason for the choice, and the last column the current state of implementation.

\begin{table}[htbp]
\centering
\small
\begin{tabularx}{\textwidth}{@{}L{3.7cm}Y L{4.4cm}@{}}
\toprule
\textbf{Method} & \textbf{ Role in the research} & \textbf{Rationale and current status} \\
\midrule
{}[System-level simulation] & [Answers RQ1 by predicting ... under conditions ...] & [Chosen over A and B because ...; model implemented and benchmarked against ...] \\
\addlinespace
{}[Convex optimisation] & [Answers RQ2 by determining ...] & [Handles the coupling between ...; formulation derived, solver tested] \\
\addlinespace
{}[Experimental validation] & [Answers RQ3 by measuring ...] & [Provides evidence that simulation cannot; rig commissioned, first data collected] \\
\bottomrule
\end{tabularx}
\caption{Methods used in the project, their role and current status}
\label{tab:methods}
\end{table}

\subsection{Data, Equipment and Experimental Setup}
% 数据从哪来、设备是什么、参数怎么定、如何保证可复现。
[Describe the data sources and their volume, the equipment and software used, the parameter
ranges and how they were chosen, and the steps taken to make the work reproducible: version
control, fixed random seeds, documented calibration procedures.]

\subsection{Validity, Reproducibility and Limitations}
The main limitations of the methodology are [limitation 1], [limitation 2] and [limitation 3].
The first is addressed by [mitigation]; the second by [mitigation]; the third is accepted,
because [reason], and its effect on the conclusions is [bounded in this way]. Where a limitation
could not be removed, the alternative approach that would be adopted is described in
Section 5.3.

% ============ 4. 已完成工作：对照原计划，讲清结果 ============
\section{Progress During the Review Period}

\subsection{Work Completed}
% 对照注册时商定的计划逐项说明；有偏差就解释，不要回避。
The plan agreed with the supervisory team at registration set out five milestones for the first
year. Table~\ref{tab:milestones} compares the plan with what has been achieved. Four of the five
milestones have been met in full; the platform milestone has been met in part, for the reason
explained in Section 4.2.

\begin{table}[htbp]
\centering
\small
\begin{tabularx}{\textwidth}{@{}L{5.4cm}L{2.2cm}L{1.9cm}Y@{}}
\toprule
\textbf{Milestone} & \textbf{Planned} & \textbf{Status} & \textbf{Comment} \\
\midrule
Literature review & M1--M4 & Completed & [Number] papers reviewed; review chapter drafted \\
\addlinespace
Method selection and training & M3--M5 & Completed & [Method] selected and benchmarked \\
\addlinespace
Platform, rig or code base & M4--M8 & Partly completed & [What is working; what is outstanding] \\
\addlinespace
Preliminary results & M8--M11 & Completed & [Dataset of N cases; results in Section 4.2] \\
\addlinespace
First paper draft & M10--M12 & Completed & Draft submitted to supervisors in [month] \\
\bottomrule
\end{tabularx}
\caption{Milestones planned at registration and their status}
\label{tab:milestones}
\end{table}

The work carried out in WP1 concentrated on [what was done], and produced [the output]. In WP2,
[what was done] was implemented and tested against [the benchmark]. The platform work required
more time than expected because [the technical reason], which is discussed further in
Section 7.1.

\subsection{Preliminary Results}
% 放 1--2 张最能说明"进展是实的"的图，并给出解读，而不是只贴图。
Figure~\ref{fig:result1} shows [what the figure shows]. [Describe the trend or comparison.] The
result is consistent with [expectation from the literature], and it suggests that [the
interpretation relevant to RQ1]. It also reveals [the unexpected feature], which will be
examined in WP2.

\begin{figure}[htbp]
\centering
% \includegraphics[width=0.75\textwidth]{image/result1.png}
\fbox{\parbox[c][5cm][c]{0.75\textwidth}{\centering Placeholder for the first key result}}
\caption{[Caption: what is plotted, under which conditions, and what the reader should notice.]}
\label{fig:result1}
\end{figure}

[If a second result is needed, present it in the same way: one paragraph of interpretation per
figure, no more.]

\subsection{Research Outputs}
The following outputs have been produced during the review period:
\begin{itemize}[leftmargin=*]
    \item Published: [authors, title, venue, year, DOI]
    \item Under review: [title, venue, submission date, status]
    \item In preparation: [title, target venue, intended submission date]
    \item Presentations and other outputs: [talks, posters, datasets, code, patents]
\end{itemize}

% ============ 5. 计划：路径、甘特图、风险、论文时间表 ============
\section{Research Plan for the Remaining Period}

\subsection{Work Packages and Timeline}
% 计划要落到"时间 + 交付物"，并且和前面的 RQ/目标对得上。
The remaining work is organised into three work packages. WP2 will [what will be done] and
deliver [the deliverable] by month [18]. WP3 will [what will be done] and deliver [the
deliverable] by month [28]. WP4 covers [integration, writing and revision] and leads to
submission in month [36]. Table~\ref{tab:plan} gives the schedule for the next twelve months in
more detail; the full project timeline is shown in Figure~\ref{fig:gantt}.

\begin{table}[htbp]
\centering
\small
\begin{tabularx}{\textwidth}{@{}L{2.2cm}Y L{4.6cm}@{}}
\toprule
\textbf{Month} & \textbf{Activities} & \textbf{Output} \\
\midrule
13--15 & [WP2: formulation, implementation and testing] & [Working method and test results] \\
\addlinespace
16--18 & [WP2: parametric study and analysis] & [Dataset and analysis] \\
\addlinespace
19--21 & [WP3: experimental campaign] & [Measurements and validation report] \\
\addlinespace
22--24 & [WP3: comparison, paper writing, Stage 2 preparation] & [Journal draft and Stage 2 report] \\
\bottomrule
\end{tabularx}
\caption{Plan for the next twelve months}
\label{tab:plan}
\end{table}

\subsection{Updated Gantt Chart}
The Gantt chart in Figure~\ref{fig:gantt} shows the whole project. Completed tasks are shown in
full; the red markers indicate the progression reviews at months 12, 24 and 30, and thesis
submission at month 36.
\begin{landscape}
\begin{figure}[htbp]
\centering
\begin{ganttchart}[
    y unit title=0.45cm,
    y unit chart=0.62cm,
    vgrid,
    title height=0.9,
    bar/.style={fill=blue!50},
    incomplete/.style={fill=green!50},
    milestone/.style={fill=red!70},
    progress label text={},
    bar height=0.62,
    group right shift=0,
    group top shift=0.7,
    group height=.3,
    group peaks width={0.2},
    milestone right shift=0,
    milestone top shift=0.7,
    x unit=0.4cm,
    canvas/.style={draw=none}
    ]{1}{36}
    \gantttitle{Research Plan Timeline (Months)}{36} \\
    \gantttitlelist{1,...,36}{1} \\

    \ganttgroup{Year 1: Foundations}{1}{12} \\
    \ganttbar[progress=100]{Literature review}{1}{4} \\
    \ganttbar[progress=100]{Method selection and training}{3}{5} \\
    \ganttbar[progress=100]{Platform and code base}{4}{8} \\
    \ganttbar[progress=100]{Preliminary results}{8}{11} \\
    \ganttbar[progress=100]{Stage 1 report and review}{11}{12} \\
    \ganttmilestone{Stage 1 review}{12} \\

    \ganttgroup{Year 2: Core investigation}{13}{24} \\
    \ganttbar{WP2: method development}{13}{18} \\
    \ganttbar{WP2: parametric study}{18}{20} \\
    \ganttbar{WP3: experimental campaign}{20}{23} \\
    \ganttbar{Stage 2 report and review}{23}{24} \\
    \ganttmilestone{Stage 2 review}{24} \\

    \ganttgroup{Year 3: Consolidation and thesis}{25}{36} \\
    \ganttbar{WP4: integration and analysis}{25}{28} \\
    \ganttbar{Thesis writing}{28}{35} \\
    \ganttmilestone{Stage 3 review}{30} \\
    \ganttbar{Final revision and submission}{35}{36} \\
    \ganttmilestone{Thesis submission}{36}
\end{ganttchart}
\caption{Updated project Gantt chart}
\label{fig:gantt}
\end{figure}
\end{landscape}

\subsection{Risks and Mitigation}
% 每条风险配"触发条件"，否则 mitigation 只是口号。
Four risks could affect the plan. They are listed in Table~\ref{tab:risk} together with the
mitigation measures and the point at which each will be reviewed. The most serious is [risk 1],
because [reason]; if it materialises, [the fallback].

\begin{table}[htbp]
\centering
\small
\begin{tabularx}{\textwidth}{@{}L{4.2cm}L{1.4cm}L{3.2cm}Y@{}}
\toprule
\textbf{Risk} & \textbf{Level} & \textbf{Mitigation} & \textbf{Contingency and trigger} \\
\midrule
{}[Experimental facility unavailable] & High & [Secure alternative access; validate simulation model] & [Simulation-only route if facility is not available by month 14] \\
\addlinespace
{}[Data collection slower than planned] & Medium & [Prioritise the most informative conditions; seek technical support] & [Reduce the parameter set if the campaign is not complete by month 22] \\
\addlinespace
{}[Key assumption fails] & Medium & [Test the assumption early in WP2] & [Adopt the alternative formulation; narrow the scope] \\
\addlinespace
{}[Researcher availability] & Low & [Document data and code continuously] & [Pause and resume work packages without loss of continuity] \\
\bottomrule
\end{tabularx}
\caption{Risks, mitigation measures and contingency plans}
\label{tab:risk}
\end{table}

\subsection{Thesis Completion Plan}
The thesis will be written alongside the research rather than in a separate writing-up period,
as required by the programme. Table~\ref{tab:thesis} sets out the milestones for completion and
submission within the registered period.

\begin{table}[htbp]
\centering
\small
\begin{tabularx}{\textwidth}{@{}Y L{3.4cm}@{}}
\toprule
\textbf{Milestone} & \textbf{Target} \\
\midrule
Empirical work completed & Month [28] \\
\addlinespace
First full draft to supervisors & Month [32] \\
\addlinespace
Revised draft after feedback & Month [34] \\
\addlinespace
Thesis submission & Month [36] \\
\addlinespace
Viva voce examination & Month [40] \\
\bottomrule
\end{tabularx}
\caption{Thesis completion plan}
\label{tab:thesis}
\end{table}

% ============ 6. 培训、交流与督导：标准 6 和 4 ============
\section{Training, Supervision and Professional Development}

\subsection{Researcher Training}
% 必修培训要有完成证据（证书、成绩单存于督导记录系统）。
The training activities completed or in progress during the review period are listed in
Table~\ref{tab:training}. The mandatory elements required before confirmation of registration
have been completed, and the certificates are held in the supervision records.

\begin{table}[H]
\centering
\small
\begin{tabularx}{\textwidth}{@{}Y L{3.0cm}L{2.2cm}@{}}
\toprule
\textbf{Activity} & \textbf{Provider} & \textbf{Date and status} \\
\midrule
{}[Research integrity and ethics] & [Provider] & [MM/YYYY; completed] \\
\addlinespace
{}[Health, safety and laboratory induction] & [Provider] & [MM/YYYY; completed] \\
\addlinespace
{}[Discipline-specific module] & [Provider] & [MM/YYYY; passed] \\
\addlinespace
{}[Technical or transferable skills course] & [Provider] & [MM/YYYY; in progress] \\
\bottomrule
\end{tabularx}
\caption{Training activities during the review period}
\label{tab:training}
\end{table}

\subsection{Academic Engagement and Communication}
% 对应标准 4：向不同受众讲清自己的研究；同行活动与科普活动都算。
Beyond formal training, I have taken part in [the weekly group meetings, the school seminar
series, the reading group], and attended [number] seminars during the year. I presented
[the work] at [the group seminar / the departmental poster session / the conference], adapting
the level of detail to the audience: the group presentation focused on [the technical detail],
whereas the poster and the three-minute presentation for [the wider audience] concentrated on
[the problem and its significance]. A plain-language description of the project was also
prepared for [the outreach event or the school newsletter].

In the coming year I plan to present at [conference], and to submit [the first journal paper]
by month [18].

\subsection{Supervision and Research Environment}
Supervisory meetings have been held [weekly / fortnightly / monthly], as required by the
programme. Each meeting followed an agreed agenda covering progress against the plan, problems
encountered, and the next steps; the notes were agreed with the supervisors and are held in the
supervision records. In addition to the scheduled meetings, I had regular informal discussions
with [the principal supervisor] on [the specific topic].

The research environment is adequate for the planned work: [the laboratory, the computing
facility, the software and the technical support] are available, and [the specific resource] was
secured in [month]. The one environment-related risk is [the risk], which is included in
Table~\ref{tab:risk}.

% ============ 7. 困难与所需支持：如实写 ============
\section{Challenges and Support Required}

\subsection{Challenges Encountered}
% 写清"困难、影响、已做的应对"，不必回避，也不要只写客观原因。
The main difficulties encountered during the year were [challenge 1], [challenge 2] and
[challenge 3]. [Challenge 1] delayed [the task] by [duration]; it was addressed by [action], at
the cost of [consequence]. [Challenge 2] affected [the work], and was mitigated by [action].
[Personal circumstances, if any, should be described here with their timing and effect, and
whether an exceptional circumstances claim was submitted.]

\subsection{Changes to the Plan}
As a result of these difficulties, [the milestone] was moved from month [x] to month [y], and
[the method or scope] was adjusted so that [the reason]. The revised plan in Section 5 already
reflects these changes, and the remaining milestones are unaffected.

\subsection{Support Required}
The following support would help to deliver the revised plan:
\begin{itemize}[leftmargin=*]
    \item Access to [the test facility] for two weeks in month 15, requested from [the
    laboratory manager].
    \item A licence for [the optimisation toolbox], requested from [the school].
    \item Travel funding for [the conference] in month 20.
\end{itemize}

% ============ 8. 结论 ============
\section{Conclusion}

This report has set out the background and rationale of the project, reviewed the literature
and identified the gap, described the methodology, and reported the progress made during the
first year. The main achievements are [the three achievements, in one sentence each]. The
preliminary results support [the interpretation] and provide a basis for the work in WP2.

The remaining work is organised into [number] work packages with defined deliverables, and the
plan leads to thesis submission within the registered period. The main risks have been
identified and mitigation measures are in place.

On this basis, I consider that the Stage 1 requirements have been met, and I am submitting this
report for confirmation of PhD registration.

% Reference functions
\addrefstosection
\printreferences

\end{document}
